% ============================================================
% The Spectral Structure of Fractional Fields
% ============================================================

\documentclass[12pt]{amsart}

\usepackage{amsmath,amssymb,amsthm}
\usepackage{booktabs}
\usepackage{microtype}
\usepackage{url}

\newtheorem{theorem}{Theorem}
\newtheorem{lemma}[theorem]{Lemma}
\newtheorem{proposition}[theorem]{Proposition}
\newtheorem{corollary}[theorem]{Corollary}
\theoremstyle{definition}
\newtheorem{definition}[theorem]{Definition}
\theoremstyle{remark}
\newtheorem{remark}[theorem]{Remark}

\title{The Spectral Structure of Fractional Fields}

\author{Alexander S.\ Petty}

\email{alexander.petty@gmail.com}
\date{July 2021 (revised April 2026)}
\subjclass[2020]{11A63, 11B83, 42A16}

\begin{document}
\begin{abstract}
For primitive-root primes ($\operatorname{ord}_p(b) = p{-}1$),
we prove that the cross-alignment matrix, reordered by
discrete-logarithm offset, is a circulant whose eigenvalues
are the Fourier transform of the digit-equality
autocorrelation. For primes with exactly two cosets and
zero within-coset autocorrelation at nonzero lags, the
spectrum has two eigenvalue levels determined by the
cross-coset match count. The multi-coset case is described
observationally.
\end{abstract}
\maketitle

% ============================================================
\section{Introduction}

The cross-alignment matrix $\mathbf{A}(p)$ introduced
in~\cite{paper6} is a symmetric, positive semidefinite
matrix on the fractions of a prime $p$ in base $b$, whose
$(i,j)$ entry is the digit-equality proportion between
$i/p$ and $j/p$. For digit-partitioning primes
($p \le b{+}1$), the matrix is the identity, every
fraction is a stranger to every other, and the eigenvalue
spectrum is trivially flat. For non-digit-partitioning
primes, the matrix is no longer the identity. Its
eigenvalue spectrum encodes the interference geometry of
the field, but the spectrum was not given a closed form
in~\cite{paper6}.

This paper identifies the eigenvalues. The main result
(Theorem~\ref{thm:circulant}) is that for primes $p$
where $b$ is a primitive root modulo $p$, the matrix
$\mathbf{A}(p)$, after a discrete-logarithm reordering of
its rows and columns, becomes a circulant matrix whose
first row is the cyclic autocorrelation of the base
repetend $1/p$, divided by the cycle length $L = p - 1$.
Standard properties of circulants~\cite{davis} then
identify the eigenvalues as the discrete Fourier
transform of the first row. The eigenvalue spectrum of
$\mathbf{A}(p)$ is a character expansion of the
repetend's autocorrelation.

For primes $p$ where $b$ is not a primitive root, the
fractions partition into $C = (p{-}1)/L$ orbits of
$\langle b \rangle$, and the matrix decomposes into
within-orbit and between-orbit blocks. We treat the
two-orbit case in detail (Theorem~\ref{thm:two-coset})
and show that under the natural sufficient conditions,
which are satisfied at $p = 13$ in base $10$, the
spectrum collapses to two eigenvalue levels determined
by an involution. The general multi-orbit case is
described observationally in~\S\ref{sec:multi-coset}.

The connection to harmonic analysis on finite groups is
explicit: the eigenvalues are inner products of the
autocorrelation function with the additive characters of
$\mathbb{Z}/L\mathbb{Z}$. This identifies the eigenvalue
spectrum as a Fourier transform on a finite cyclic group,
which is the natural starting point for the harmonic
analysis pursued in the next several papers.

% ============================================================
\section{The Circulant Theorem}

\begin{definition}
The \emph{cyclic autocorrelation} of the base repetend
$(d_0, d_1, \ldots, d_{L-1})$ of $1/p$ is the sequence
\[
R(\ell) = |\{j \in \{0, \ldots, L{-}1\} : d_j = d_{(j+\ell)
\bmod L}\}|,
\qquad \ell = 0, 1, \ldots, L{-}1.
\]
\end{definition}

Note that $R(0) = L$ (every position matches itself) and
$R(\ell) \ge 0$ for all $\ell$.

\begin{theorem}\label{thm:circulant}
Let $p$ be a prime for which $b$ is a primitive root
($\operatorname{ord}_p(b) = p{-}1$). Then the cross-alignment
matrix $\mathbf{A}(p)$, reordered by discrete-logarithm
offset, is a circulant matrix with first row
\[
\left(\frac{R(0)}{L},\; \frac{R(1)}{L},\;
\frac{R(2)}{L},\; \ldots,\; \frac{R(L{-}1)}{L}\right).
\]
Its eigenvalues are given by the discrete Fourier
transform~\cite{davis}:
\[
\lambda_j = \frac{1}{L} \sum_{\ell=0}^{L-1}
R(\ell)\, \omega^{j\ell},
\qquad j = 0, 1, \ldots, L{-}1,
\]
where $\omega = e^{2\pi i / L}$.
\end{theorem}

\begin{proof}
When $b$ is a primitive root modulo $p$, every fraction $k/p$ is a
cyclic rotation of $1/p$ in the sense that the repetend of $k/p$ in
base $b$ is the repetend of $1/p$ shifted by $\log_b k$ positions
modulo $L$ (where $\log_b$ denotes the discrete logarithm with
respect to $b$ modulo $p$). Ordered by rotation offset, the
fractions are $1/p,\ b/p,\ b^2/p,\ \ldots,\ b^{L-1}/p$ (each
reduced modulo $p$).

In this ordering, the $(i, j)$ entry of $\mathbf{A}(p)$ is the
digit-equality proportion between $b^i/p$ and $b^j/p$, which is
the proportion of positions in the repetend of $1/p$ where the
digit at position $\ell$ agrees with the digit at position
$\ell + (j - i) \bmod L$. By the definition of $R$, this proportion
is $R(j - i \bmod L)/L$. The $(i,j)$ entry depends only on the
difference $j - i$ modulo $L$, so the matrix is circulant
with first row $(R(0)/L,\ R(1)/L,\ \ldots,\ R(L-1)/L)$.

The eigenvalues of a circulant are the discrete Fourier transform
of its first row~\cite{davis}, which yields the stated formula.
\end{proof}

\begin{corollary}\label{cor:dp-flat}
If $p$ is digit-partitioning ($p \le b{+}1$), every eigenvalue
of $\mathbf{A}(p)$ equals $1$.
\end{corollary}

\begin{proof}
By~\cite{paper6}, $\mathbf{A}(p) = I_{p-1}$ for digit-partitioning
$p$, so the spectrum is the constant sequence $1$. For
primitive-root digit-partitioning primes, the same conclusion
follows from Theorem~\ref{thm:circulant}: the digit-partitioning
condition forces $R(\ell) = 0$ for $\ell \ne 0$ (no two distinct
positions of the repetend share a digit), and the DFT of
$(L, 0, 0, \ldots, 0)/L$ is the constant sequence $1$.
\end{proof}

% ============================================================
\section{Worked Example: $p = 7$ in Base $10$}\label{sec:p7}

In base $10$, the prime $p = 7$ is digit-partitioning
($p \le b{+}1 = 11$) and primitive-root
($\operatorname{ord}_7(10) = 6$). The repetend of $1/7$ is
$142857$, of length $L = 6$. By Theorem~\ref{thm:circulant},
$\mathbf{A}(7)$ in the discrete-logarithm ordering is a
$6 \times 6$ circulant whose first row is
$(R(0)/6, R(1)/6, \ldots, R(5)/6)$.

The autocorrelation values are computed by direct comparison
of $142857$ with its rotations. At lag $\ell = 1$ the rotated
sequence is $428571$, and the position-wise comparison
$(1{:}4,\ 4{:}2,\ 2{:}8,\ 8{:}5,\ 5{:}7,\ 7{:}1)$ yields zero
matches. The same calculation at every lag $\ell \in
\{1, 2, 3, 4, 5\}$ gives zero matches. Hence
$R(0) = 6$ and $R(\ell) = 0$ for $\ell \ne 0$, and the first
row of the circulant is $(1, 0, 0, 0, 0, 0)$.

The DFT of this row is the constant sequence
$(1, 1, 1, 1, 1, 1)$. All six eigenvalues of $\mathbf{A}(7)$
equal $1$, in agreement with the digit-partitioning case
of Corollary~\ref{cor:dp-flat}.

% ============================================================
\section{The Two-Orbit Case}

When $\operatorname{ord}_p(b) = L < p{-}1$, the fractions
partition into $C = (p{-}1)/L$ orbits of $\langle b \rangle$
in $(\mathbb{Z}/p\mathbb{Z})^*$. For $C = 2$, the
cross-alignment matrix has additional structure that we
make precise below.

\begin{theorem}\label{thm:two-coset}
Let $p$ be a prime with exactly two orbits of $\langle b
\rangle$ in $(\mathbb{Z}/p\mathbb{Z})^*$, each of size
$L = (p{-}1)/2$. Suppose
\begin{enumerate}
\item the within-orbit autocorrelation
$R_C(\ell)$ (the digit-equality count between rotations
within the same orbit) satisfies $R_C(\ell) = 0$ for
$\ell \ne 0$, and
\item each fraction in one orbit matches exactly one
fraction in the other orbit at exactly $m$ positions, with
all other cross-orbit pairs matching at zero positions.
\end{enumerate}
Then the eigenvalues of $\mathbf{A}(p)$ are
\[
1 + \frac{m}{L} \quad \text{and} \quad 1 - \frac{m}{L},
\]
each with multiplicity $L = (p{-}1)/2$.
\end{theorem}

\begin{proof}
Under these conditions, every diagonal entry of $\mathbf{A}$
is $1$, every within-orbit off-diagonal entry is $0$ (from
condition (1)), and every cross-orbit off-diagonal entry is
either $m/L$ (for the unique partner) or $0$ (otherwise).
We can therefore write
\[
\mathbf{A}(p) = \mathbf{I} + \frac{m}{L}\, \mathbf{P},
\]
where $\mathbf{P}$ is the permutation matrix that pairs each
fraction with its unique cross-orbit partner. Since the
pairing is symmetric, $\mathbf{P}$ is an involution
($\mathbf{P}^2 = \mathbf{I}$), so its eigenvalues are
$\pm 1$, each with multiplicity equal to half its dimension,
i.e., multiplicity $L$ each. The eigenvalues of $\mathbf{A}$
follow as $1 \pm m/L$.
\end{proof}

\begin{proposition}\label{prop:p13-verify}
The hypotheses of Theorem~\ref{thm:two-coset} are satisfied
at $p = 13$ in base $10$, with $m = 2$ and $L = 6$.
The eigenvalues of $\mathbf{A}(13)$ are $4/3$ and $2/3$,
each with multiplicity $6$.
\end{proposition}

\begin{proof}
The order of $10$ modulo $13$ is $6$, so $L = 6$ and
$C = 2$. The two orbits of $\langle 10 \rangle$ in
$(\mathbb{Z}/13\mathbb{Z})^*$ are
$\{1, 10, 9, 12, 3, 4\}$ and $\{2, 7, 5, 11, 6, 8\}$. The
repetends within the first orbit are the rotations of
$076923$, and within the second orbit they are the rotations
of $153846$.

Within-orbit comparison: comparing $076923$ with each of its
rotations $769230$, $692307$, $923076$, $230769$, $307692$
position by position gives zero matches in every case.
Hence $R_C(\ell) = 0$ for $\ell \ne 0$, verifying
condition (1).

Cross-orbit comparison: comparing $1/13 = 076923$ with each
member of the second orbit gives the following position-wise
match counts:
\begin{center}
\begin{tabular}{ll}
\toprule
partner & matches \\
\midrule
$2/13 = 153846$ & $0$ \\
$5/13 = 384615$ & $0$ \\
$6/13 = 461538$ & $0$ \\
$7/13 = 538461$ & $0$ \\
$8/13 = 615384$ & $0$ \\
$11/13 = 846153$ & $2$ \\
\bottomrule
\end{tabular}
\end{center}
The unique partner of $1/13$ in the other orbit is $11/13$,
matching at exactly $2$ positions (positions $3$ and $6$).
The same pattern holds for every other fraction in the first
orbit, by replaying the calculation: each fraction has a
unique cross-orbit partner matching at exactly $2$ positions
and matches every other cross-orbit fraction at $0$
positions. This verifies condition (2) with $m = 2$.

By Theorem~\ref{thm:two-coset}, the eigenvalues of
$\mathbf{A}(13)$ are $1 + 2/6 = 4/3$ and $1 - 2/6 = 2/3$,
each with multiplicity $6$.
\end{proof}

Whether the hypotheses of Theorem~\ref{thm:two-coset}
hold for every two-orbit prime, or only for a subfamily
that includes $p = 13$, is not determined here.

% ============================================================
\section{The General Multi-Orbit Case}\label{sec:multi-coset}

For primes with $C = (p{-}1)/L > 2$ orbits, the cross-alignment
matrix has a block decomposition that we describe but do not
prove in full generality.

\begin{proposition}\label{prop:block-decomp}
For any prime $p$ with $C$ orbits of $\langle b \rangle$
in $(\mathbb{Z}/p\mathbb{Z})^*$, each of size $L$, the
cross-alignment matrix admits the decomposition
\[
\mathbf{A}(p) = \mathbf{W}(p) + \mathbf{X}(p),
\]
where $\mathbf{W}(p)$ is the block-diagonal within-orbit
contribution (a direct sum of $C$ circulant $L \times L$
matrices, one per orbit, each governed by the within-orbit
autocorrelation) and $\mathbf{X}(p)$ is the
block-off-diagonal between-orbit contribution (whose
$(i,j)$ block records the cross-orbit alignments between
orbit $i$ and orbit $j$).
\end{proposition}

\begin{proof}
The block partition follows from grouping the rows and
columns of $\mathbf{A}(p)$ by orbit. Within each orbit
of size $L$, the rotation argument of
Theorem~\ref{thm:circulant} applies verbatim, since
fractions within the same orbit are rotations of one
another, so each diagonal block is a circulant whose first
row is $R_C(\ell)/L$ for that orbit.
The cross-orbit blocks contain the digit-equality
proportions between fractions of different orbits and have
no cyclic structure in general.
\end{proof}

The within-orbit eigenvalues are the DFT of $R_C(\ell)/L$
for each orbit, by Theorem~\ref{thm:circulant} applied
within the orbit. The cross-orbit blocks add a perturbation
whose eigenvalues depend on the joint structure of all
$C$ orbits.

Computation for all primes $p \le 37$ in base $10$ shows
that the spectrum of $\mathbf{A}(p)$ is determined by the
within-orbit autocorrelations $R_C(\ell)$ together with the
cross-orbit alignment data, in the sense that two primes
with the same orbit structure and the same alignment data
have the same spectrum. A closed-form expression for the
multi-orbit eigenvalues in terms of these data is not
given here.

% ============================================================
\section{Connection to Harmonic Analysis}

The DFT that produces the eigenvalues is a sum over characters
of $\mathbb{Z}/L\mathbb{Z}$:
\[
\lambda_j = \frac{1}{L} \sum_{\ell=0}^{L-1}
R(\ell)\, \chi_j(\ell),
\]
where $\chi_j(\ell) = e^{2\pi i j\ell/L}$ is the $j$-th character.
The eigenvalue $\lambda_j$ is therefore the inner product of the
autocorrelation function with the character $\chi_j$.

This identifies the eigenvalue spectrum of the cross-alignment
matrix as a \emph{character expansion} of the repetend's
autocorrelation, in the sense of harmonic analysis on finite
abelian groups~\cite{terras}. The digit-partitioning property corresponds to a
flat character expansion (all coefficients equal). Non-DP primes
exhibit character-dependent structure.

For primitive-root primes, where $L = p - 1$, the additive
characters of $\mathbb{Z}/L\mathbb{Z}$ used here are not the
Dirichlet characters of $(\mathbb{Z}/p\mathbb{Z})^*$, but the
two character groups are isomorphic via the discrete logarithm:
the map $\log_b : (\mathbb{Z}/p\mathbb{Z})^* \to \mathbb{Z}/L\mathbb{Z}$
intertwines a multiplicative character of
$(\mathbb{Z}/p\mathbb{Z})^*$ with an additive character of
$\mathbb{Z}/L\mathbb{Z}$. The eigenvalues of $\mathbf{A}(p)$ can
therefore be re-expressed as evaluations of Dirichlet characters
modulo $p$ on the autocorrelation pulled back through the
discrete logarithm. The Dirichlet character framework~\cite{iwaniec}
provides the analytic machinery for relating these character sums
to $L$-functions, and is the natural setting in which the next
several papers in the chain are developed.

% ============================================================
\section{Remarks}

\subsection*{The autocorrelation in the primitive-root case}

For primitive-root primes, the cyclic autocorrelation
$R(\ell)$ determines the full eigenvalue spectrum
(Theorem~\ref{thm:circulant}). For multi-coset primes,
the within-coset autocorrelation alone does not suffice:
the between-coset cross-correlation data is also needed
(as seen at $p = 13$ in base~$10$, where $R(\ell) = 0$
for $\ell \ne 0$ yet the spectrum is not flat).

\subsection*{From spectra to $L$-functions}

The expression $\lambda_j = L^{-1} \sum R(\ell) \chi_j(\ell)$
resembles a Gauss-type sum. If the autocorrelation $R(\ell)$ can
be expressed in terms of the Legendre symbol or other
multiplicative characters, the eigenvalues become evaluations of
Dirichlet $L$-functions at specific points. Whether this connection
yields non-trivial analytic information about the fractional field
is open.

% ============================================================
\begin{thebibliography}{9}

\bibitem{paper6}
A.~S.~Petty,
The cross-alignment matrix,
research note, 2021.

\bibitem{davis}
P.~J.~Davis,
\emph{Circulant Matrices},
2nd ed., Chelsea, 1994.

\bibitem{terras}
A.~Terras,
\emph{Fourier Analysis on Finite Groups and Applications},
Cambridge University Press, 1999.

\bibitem{iwaniec}
H.~Iwaniec and E.~Kowalski,
\emph{Analytic Number Theory},
Amer.\ Math.\ Soc., 2004.

\end{thebibliography}

\end{document}
